\section{ML-DSA: Post-Quantum Signatures}


\subsection{FIPS 204 Overview}

ML-DSA is a lattice-based signature scheme built on the Module Learning With Errors and Module Short Integer Solution (MSIS) problems.

\begin{definition}[ML-DSA Parameters]
\begin{itemize}
\item \textbf{ML-DSA-65}: NIST security level 3. Default for transaction signing.
\item \textbf{ML-DSA-87}: NIST security level 5. Required for long-term credential signing and governance actions.
\end{itemize}
\end{definition}

\subsection{Signature Sizes}

\begin{table}[h]
\centering
\begin{tabular}{lccc}
\toprule
\textbf{Scheme} & \textbf{Public Key} & \textbf{Signature} & \textbf{Security Level} \\
\midrule
ML-DSA-65 & 1,952 bytes & 3,293 bytes & 3 \\
ML-DSA-87 & 2,592 bytes & 4,595 bytes & 5 \\
Ed25519 & 32 bytes & 64 bytes & $\approx$1 (classical) \\
BLS12-381 & 48 bytes & 96 bytes & $\approx$1 (classical) \\
\bottomrule
\end{tabular}
\caption{Signature sizes (ML-DSA vs classical)}
\label{tab:dsa-sizes}
\end{table}

\subsection{Threshold ML-DSA}

Threshold ML-DSA signing uses a commit-reveal structure: each signer $v_i$ commits $w_i = A \cdot y_i$ (masking vector), the challenge $c = H(\text{msg} \| w_1 \| \cdots \| w_t)$ is computed after all commitments, each $v_i$ responds with $z_i = y_i + c \cdot s_i$ (with rejection sampling), and \KeyVM{} aggregates $z = \sum_{i \in S} \lambda_i \cdot z_i$.

\begin{theorem}[Threshold ML-DSA Security]
Under MLWE and MSIS hardness, the threshold ML-DSA scheme is EUF-CMA secure when the adversary controls fewer than $t$ signers.
\end{theorem}

